\begin{Exercise}[label={ex:voti}]
    Si scriva un programma che gestisca un database di voti universitari salvati su un file "registro.txt".
    Il programma deve gestire due modilità operative distinte:
    \begin{itemize}
        \item Modalità aggiunta (-add): aggiunge un nuovo studente al registro. Richiede tre parametri: nome, cognome e voto. 
        \item Modalità ricerca (-find): cerca e stampa a video tutti gli studenti che hanno un voto superiore o uguale ad una soglia specificata.
    \end{itemize}

    Si scriva una funzione \texttt{aggiungi\_studente()} che aggiunge un nuovo studente al file "registro.txt" 
    e una funzione \texttt{cerca\_studenti()} che apre e cerca all'interno di "registro.txt" 
    tutti gli studenti con voto superiore o uguale alla soglia specificata. Si gestiscano le due modalità come
    argomenti da riga di comando e si aggiunga la possibilità di utilizzare un file di registro personalizzato usando come flag "-f".

    \textbf{Esempio:}
    \begin{lstlisting}[style=console]
        # Usa il file di default, aggiunge Mario Rossi con 28
        ./programma -add Mario Rossi 28

        # Usa un file diverso, cerca voti >= 25
        ./programma -f esame_c.txt -find 25
    \end{lstlisting}
    
\end{Exercise}

\begin{Answer}[ref={ex:voti}]
  	\lstinputlisting[style=cstylesmall]{src/voti.c}
\end{Answer}